\chapter{収束定理}

\paragraph{本章の目標}
本章ではLebesgue積分論の中心である収束定理を学ぶ．
特に
\[
\text{単調収束定理},\qquad
\text{Fatouの補題},\qquad
\text{優収束定理},\qquad
\text{有界収束定理}
\]
を証明し，
さらに
\[
\text{項別積分},\qquad \text{積分記号下の微分}
\]
へ応用する．
これらは，確率論では期待値と極限の交換を，
関数解析では級数・極限・積分の交換を支える基本定理である．

\section{動機づけ}

第1章で見たように，Riemann積分では極限と積分の交換は繊細である．
例えば，Riemann可積分関数列が点ごとに収束しても，
極限関数はRiemann可積分とは限らない．
Dirichlet関数の例はその典型であった．

Lebesgue積分の強みは，極限の交換を可能にする十分条件が明確である点にある．
実際，現代数学では
\[
f_n \to f
\]
という極限に対して
\[
\int f_n\,dm \to \int f\,dm
\]
を言いたい場面が非常に多い．
本章では，そのための代表的な4つの定理を学ぶ．

\section{単調収束定理}

\begin{theorem}[単調収束定理]
$E \subset \RR^d$ を可測集合とし，
$f_n:E \to [0,\infty]$ を非負可測関数列とする．
もし
\[
f_1\le f_2\le \cdots,
\qquad
f_n(x)\to f(x)\quad (x \in E)
\]
ならば
\[
\int_E f_n\,dm \uparrow \int_E f\,dm
\]
すなわち
\[
\lim_{n\to\infty}\int_E f_n\,dm=\int_E f\,dm
\]
が成り立つ．
\end{theorem}
\begin{proof}
まず $f_n \le f$ だから単調性より
\[
\int_E f_n\,dm\le \int_E f\,dm
\]
であり，列 $\left\{\int_E f_n\,dm\right\}$ は単調増加で上に有界である．
したがって極限
\[
L:=\lim_{n\to\infty}\int_E f_n\,dm
\]
が存在して
\[
L\le \int_E f\,dm
\]
である．

逆向きの不等式を示す．
$s \le f$ を満たす任意の非負単関数
\[
s=\sum_{j=1}^m a_j1_{A_j}
\]
を取る．
ここで $a_j\ge 0$，$A_j$ は互いに素な可測集合とする．
任意の $0<\alpha<1$ に対して
\[
B_{j,n}:=A_j \cap \{x \in E \mid f_n(x)\ge \alpha a_j\}
\]
とおく．
$f_n(x)\uparrow f(x)\ge a_j$ が $A_j$ 上で成り立つので
\[
B_{j,1}\subset B_{j,2}\subset \cdots,
\qquad
\bigcup_{n=1}^{\infty}B_{j,n}=A_j
\]
である．

そこで
\[
u_n:=\sum_{j=1}^m \alpha a_j 1_{B_{j,n}}
\]
とおくと，$u_n$ は非負単関数で
\[
u_n\le f_n
\]
が成り立つ．
したがって
\[
\int_E u_n\,dm\le \int_E f_n\,dm\le L.
\]
一方，測度の下からの連続性より
\[
m(B_{j,n})\to m(A_j)
\]
だから
\[
\int_E u_n\,dm
=
\sum_{j=1}^m \alpha a_j m(B_{j,n})
\to
\sum_{j=1}^m \alpha a_j m(A_j)
=
\alpha \int_E s\,dm.
\]
よって
\[
\alpha \int_E s\,dm\le L.
\]
$\alpha \uparrow 1$ とすれば
\[
\int_E s\,dm\le L.
\]
$s \le f$ を満たす単関数 $s$ は任意だから，
$f$ の積分の定義より
\[
\int_E f\,dm\le L.
\]
したがって
\[
\int_E f\,dm=L
\]
である．
\end{proof}

\begin{corollary}[非負級数の項別積分]
$f_n:E \to [0,\infty]$ を非負可測関数列とする．
このとき
\[
\int_E \left(\sum_{n=1}^{\infty}f_n\right)\,dm
=
\sum_{n=1}^{\infty}\int_E f_n\,dm
\]
が成り立つ．
\end{corollary}
\begin{proof}
部分和
\[
S_N:=\sum_{n=1}^N f_n
\]
を考えると
\[
S_1\le S_2\le \cdots,
\qquad
S_N \uparrow \sum_{n=1}^{\infty}f_n.
\]
したがって単調収束定理より
\[
\int_E \left(\sum_{n=1}^{\infty}f_n\right)\,dm
=
\lim_{N\to\infty}\int_E S_N\,dm.
\]
一方，有限和の線形性より
\[
\int_E S_N\,dm
=
\sum_{n=1}^N \int_E f_n\,dm.
\]
よって
\[
\int_E \left(\sum_{n=1}^{\infty}f_n\right)\,dm
=
\lim_{N\to\infty}\sum_{n=1}^N \int_E f_n\,dm
=
\sum_{n=1}^{\infty}\int_E f_n\,dm.
\]
\end{proof}

\section{Fatouの補題}

\begin{theorem}[Fatouの補題]
$f_n:E \to [0,\infty]$ を非負可測関数列とする．
このとき
\[
\int_E \liminf_{n\to\infty} f_n\,dm
\le
\liminf_{n\to\infty}\int_E f_n\,dm
\]
が成り立つ．
\end{theorem}
\begin{proof}
各 $n$ に対して
\[
g_n:=\inf_{k\ge n}f_k
\]
とおく．
すると $g_n$ は非負可測で，
\[
g_1\le g_2\le \cdots,
\qquad
g_n \uparrow \liminf_{n\to\infty} f_n.
\]
したがって単調収束定理より
\[
\int_E \liminf_{n\to\infty} f_n\,dm
=
\lim_{n\to\infty}\int_E g_n\,dm.
\]
しかも各 $n$ と各 $k\ge n$ に対して
\[
g_n\le f_k
\]
だから単調性より
\[
\int_E g_n\,dm\le \int_E f_k\,dm.
\]
ゆえに
\[
\int_E g_n\,dm\le \inf_{k\ge n}\int_E f_k\,dm.
\]
両辺で $n\to\infty$ とすると
\[
\int_E \liminf_{n\to\infty} f_n\,dm
\le
\lim_{n\to\infty}\inf_{k\ge n}\int_E f_k\,dm
=
\liminf_{n\to\infty}\int_E f_n\,dm.
\]
\end{proof}

\begin{remark}
Fatouの補題は，
「極限の積分は積分の極限以下とは限らないが，
少なくとも \(\liminf\) については不等式が成り立つ」
という定理である．
優収束定理の証明でも重要な役割を果たす．
\end{remark}

\section{優収束定理}

\begin{theorem}[優収束定理]
$f_n:E \to \eRR$ を可測関数列とし，
$g:E \to [0,\infty]$ を可積分な関数とする．
もし
\[
f_n(x)\to f(x) \quad \text{a.e.\ on } E
\]
かつ
\[
|f_n(x)|\le g(x) \quad \text{a.e.\ on } E
\]
がすべての $n$ について成り立つならば，
$f$ は可積分であり，
\[
\lim_{n\to\infty}\int_E |f_n-f|\,dm=0
\]
したがって
\[
\lim_{n\to\infty}\int_E f_n\,dm=\int_E f\,dm
\]
が成り立つ．
\end{theorem}
\begin{proof}
まず極限をとれば a.e.\ で
\[
|f|\le g
\]
である．
したがって第8章の優関数による判定より $f$ は可積分である．

次に
\[
h_n:=|f_n-f|
\]
とおくと，$h_n$ は非負可測で
\[
h_n\to 0 \quad \text{a.e.\ on } E
\]
かつ
\[
0\le h_n\le |f_n|+|f|\le 2g \quad \text{a.e.\ on } E
\]
が成り立つ．
ここで
\[
2g-h_n \ge 0
\]
なので Fatou の補題を $\{2g-h_n\}$ に適用すると
\[
\int_E \liminf_{n\to\infty}(2g-h_n)\,dm
\le
\liminf_{n\to\infty}\int_E (2g-h_n)\,dm.
\]
左辺の integrand は a.e.\ で $2g$ に収束するから
\[
2\int_E g\,dm
\le
\liminf_{n\to\infty}\left(2\int_E g\,dm-\int_E h_n\,dm\right).
\]
すなわち
\[
2\int_E g\,dm
\le
2\int_E g\,dm-\limsup_{n\to\infty}\int_E h_n\,dm.
\]
よって
\[
\limsup_{n\to\infty}\int_E h_n\,dm\le 0.
\]
各 $\int_E h_n\,dm$ は非負だから
\[
\lim_{n\to\infty}\int_E h_n\,dm=0.
\]
すなわち
\[
\int_E |f_n-f|\,dm \to 0.
\]

最後に
\[
\left|\int_E f_n\,dm-\int_E f\,dm\right|
=
\left|\int_E (f_n-f)\,dm\right|
\le
\int_E |f_n-f|\,dm
\to 0
\]
より
\[
\int_E f_n\,dm\to \int_E f\,dm
\]
を得る．
\end{proof}

\begin{corollary}[有界収束定理]
$m(E)<\infty$ とし，
$f_n:E \to \eRR$ を可測関数列とする．
もし
\[
f_n(x)\to f(x)\quad \text{a.e.\ on } E
\]
かつある定数 $M>0$ が存在して
\[
|f_n(x)|\le M \quad \text{a.e.\ on } E
\]
がすべての $n$ について成り立つならば，
$f$ は可積分で
\[
\int_E f_n\,dm\to \int_E f\,dm
\]
が成り立つ．
\end{corollary}
\begin{proof}
\[
g:=M1_E
\]
とおくと
\[
\int_E g\,dm=M\,m(E)<\infty
\]
だから $g$ は可積分である．
しかも a.e.\ で
\[
|f_n|\le g
\]
であるから，優収束定理を適用すればよい．
\end{proof}

\section{項別積分}

\begin{theorem}[優収束定理による項別積分]
$f_n:E \to \eRR$ を可積分関数列とする．
さらに，ある可積分関数 $g:E \to [0,\infty]$ が存在して
\[
\sum_{n=1}^{\infty}|f_n(x)|\le g(x) \quad \text{a.e.\ on } E
\]
を満たすとする．
このとき級数
\[
\sum_{n=1}^{\infty}f_n(x)
\]
は a.e.\ で絶対収束し，その和を $f$ と書けば $f$ は可積分で，
\[
\int_E f\,dm
=
\sum_{n=1}^{\infty}\int_E f_n\,dm
\]
が成り立つ．
\end{theorem}
\begin{proof}
部分和
\[
S_N:=\sum_{n=1}^N f_n
\]
を考える．
a.e.\ で
\[
\sum_{n=1}^{\infty}|f_n(x)|\le g(x)<\infty
\]
だから，数列 $\sum f_n(x)$ は a.e.\ で絶対収束する．
その和を $f(x)$ とおけば
\[
S_N(x)\to f(x) \quad \text{a.e.}
\]
である．

さらに三角不等式より
\[
|S_N(x)|
\le
\sum_{n=1}^N |f_n(x)|
\le
\sum_{n=1}^{\infty}|f_n(x)|
\le
g(x)
\quad \text{a.e.}
\]
だから，優収束定理より
\[
\int_E S_N\,dm\to \int_E f\,dm.
\]
一方，有限和の線形性より
\[
\int_E S_N\,dm
=
\sum_{n=1}^N \int_E f_n\,dm.
\]
したがって
\[
\int_E f\,dm
=
\lim_{N\to\infty}\sum_{n=1}^N \int_E f_n\,dm
=
\sum_{n=1}^{\infty}\int_E f_n\,dm.
\]
\end{proof}

\begin{remark}
非負項級数なら単調収束定理で，
符号が混ざる級数でも絶対値で支配されていれば優収束定理で，
項別積分が正当化される．
これが「和と積分の交換」の基本形である．
\end{remark}

\section{積分記号下の微分}

\begin{theorem}[積分記号下の微分]
$I \subset \RR$ を開区間とし，$E \subset \RR^d$ を可測集合とする．
関数
\[
F:I \times E \to \RR
\]
が次を満たすとする．
\begin{enumerate}
    \item ほとんどすべての $x \in E$ に対して，$t \mapsto F(t,x)$ は $I$ 上で $C^1$ 級である
    \item 各 $t \in I$ に対して，$x \mapsto F(t,x)$ と $x \mapsto \partial_tF(t,x)$ は可測である
    \item ある $t_0 \in I$ について $F(t_0,\cdot)$ は可積分である
    \item ある可積分関数 $g:E \to [0,\infty]$ が存在して，
    \[
    |\partial_tF(t,x)|\le g(x)
    \]
    がすべての $t \in I$ と a.e.\ の $x \in E$ で成り立つ
\end{enumerate}
このとき
\[
G(t):=\int_E F(t,x)\,dm(x)
\]
は $I$ 上でよく定義されて微分可能であり，
\[
G'(t)=\int_E \partial_tF(t,x)\,dm(x)
\]
が成り立つ．
\end{theorem}
\begin{proof}
まず $G(t)$ が well-defined であることを示す．
任意の $t \in I$ と a.e.\ の $x \in E$ に対して，
1変数の微分積分学の基本定理より
\[
F(t,x)-F(t_0,x)=\int_{t_0}^t \partial_sF(s,x)\,ds
\]
である．
したがって
\[
|F(t,x)|
\le
|F(t_0,x)|+\int_{t_0}^t |\partial_sF(s,x)|\,ds
\le
|F(t_0,x)|+|t-t_0|\,g(x).
\]
右辺は可積分だから，$F(t,\cdot)$ も可積分である．
よって $G(t)$ はよく定義される．

次に $t \in I$ を固定し，
$h$ を十分小さく取って $t+h \in I$ とする．
差商
\[
Q_h(x):=\frac{F(t+h,x)-F(t,x)}{h}
\]
を考える．
仮定 (1) と平均値の定理より，
a.e.\ の $x$ に対してある $\theta=\theta(x,h)\in(0,1)$ が存在して
\[
Q_h(x)=\partial_tF(t+\theta h,x)
\]
となる．
したがって仮定 (4) より
\[
|Q_h(x)|\le g(x)
\]
が a.e.\ で成り立つ．
また $t \mapsto F(t,x)$ の微分可能性から
\[
Q_h(x)\to \partial_tF(t,x)
\quad
\text{a.e.\ on } E
\]
である．
よって優収束定理を $Q_h$ に適用できて
\[
\lim_{h\to 0}\int_E Q_h(x)\,dm(x)
=
\int_E \partial_tF(t,x)\,dm(x)
\]
を得る．
左辺は
\[
\int_E Q_h(x)\,dm(x)
=
\frac{1}{h}
\left(
\int_E F(t+h,x)\,dm(x)-\int_E F(t,x)\,dm(x)
\right)
=
\frac{G(t+h)-G(t)}{h}
\]
だから
\[
G'(t)=\int_E \partial_tF(t,x)\,dm(x)
\]
が従う．
\end{proof}

\begin{example}
$a>0$ を固定し，
\[
F(t,x):=e^{-tx}
\qquad
((t,x)\in (a,\infty)\times[0,\infty))
\]
を考える．
このとき
\[
G(t):=\int_0^{\infty}e^{-tx}\,dx
\]
は微分可能で，
\[
G'(t)=\int_0^{\infty}(-x)e^{-tx}\,dx
\]
が成り立つ．
\end{example}
\begin{proof}
各 $x\ge 0$ に対して $t\mapsto e^{-tx}$ は $C^1$ 級で，
\[
\partial_tF(t,x)=-xe^{-tx}.
\]
さらに $t\ge a$ なら
\[
|\partial_tF(t,x)|
=
xe^{-tx}
\le
xe^{-ax}
.
\]
右辺は
\[
\int_0^{\infty}xe^{-ax}\,dx<\infty
\]
を満たすので可積分である．
したがって前定理の仮定が成り立ち，
\[
G'(t)=\int_0^{\infty}(-x)e^{-tx}\,dx
\]
を得る．
\end{proof}

\section{解釈とまとめ}

本章の主役は，極限と積分の交換である．
結論だけを並べると次のようになる．
\begin{enumerate}
    \item $f_n \uparrow f$ なら
    \[
    \int f_n \to \int f
    \]
    \item 一般の非負列では
    \[
    \int \liminf f_n \le \liminf \int f_n
    \]
    \item $|f_n|\le g \in L^1$ なら
    \[
    \int f_n \to \int f
    \]
    \item 有限測度空間では，一様有界性だけで優関数を作れる
\end{enumerate}

この4つの定理により，
\[
\lim,\qquad \sum,\qquad \frac{d}{dt},\qquad \int
\]
の交換が厳密に扱えるようになった．
これがLebesgue積分論の最重要部分であり，
確率論・偏微分方程式・Fourier解析・関数解析へ進むための基礎になる．
